\cotefile{136}

\scanpage{33}

donc $C_{l,i+j}$ et $C_{l+j,i}$ $\not\equiv 0\ (p)$. C'est le contenu du

\uline{Lemme 2} Soient $i,j\geq 0$. Alors il existe $l\geq 0$ tel que \deleted{\unsure{$C_{i+j,l}$}} $C_{i+j,l}$ et $C_{i,j+l}$ soient $\not\equiv 0\ (p)$.

\deleted{\uline{Dém.} Considérons \unsure{les|des} développements $p$-adiques de $i$ et $j$

\[
\begin{aligned}
i&=i_0+i_1p+\cdots+i_rp^r\\
j&=j_0+j_1p+\cdots+j_rp^r
\end{aligned}
\]

alors le développement $p$-adique de \inserted{$\sigma=$}$i+j$ est \inserted{$\deleted{\unsure{i+}}j=\sigma=\sigma_0+\sigma_1p+\cdots+\sigma_{r+1}p^{r+1}$} donné par le procédé \unsure{vraiment|usuel} bien connu, où on détermine par récurrence des couples d'entiers $\sigma_\alpha$, \deleted{\unsure{$\rho_\alpha$}} $(0\leq\alpha\leq r+1)$ \inserted{et $\rho_\alpha$ $(-1\leq\alpha\leq r)$}

-- les $\rho_\alpha$ sont les ``retenues'', égales soit à $0$ ou $1$ -- en posant \deleted{\illegible{1}} \textcircled{a} $\rho_{-1}=0$, puis par récurrence par

\[
\underset{\downarrow}{\text{\textcircled{$0\leq\alpha\leq r\deleted{\unsure{+}}$}}}
\quad\text{\textcircled{b}}\quad i_\alpha+j_\alpha+\rho_{\alpha-1}=\sigma_\alpha+p\rho_\alpha
\qquad 0\leq\sigma_\alpha,\unsure{p|\rho}_\alpha\leq p-1
\]

celà détermine $\sigma_\alpha$ et $\rho_\alpha$

\[
\text{enfin}\quad \text{\textcircled{c}}\qquad \sigma_{\alpha+1}=\rho_{\unsure{r|\alpha}}
\]

Posons

\[
\begin{aligned}
l_\alpha&=(p-1)-\sigma_\alpha\qquad 0\leq\alpha\leq r+1\\
l&=\sum_{0}^{r+1}l_\alpha p^\alpha
\end{aligned}
\]

\[
\text{d'où}\qquad \sigma+l=\sum_0^{r+1}(p-1)p^\alpha=p^{r+2}-1\qquad \text{i.e.}\quad l=p^{r+2}-1-\sigma
\]

On a alors, grâce à (b)

\[
\text{\textcircled{$a'$}}\qquad \rho_{-1}=0
\]

\[
\text{\textcircled{$b'$}}\quad\text{Si }0\leq\alpha\leq r,\quad
(i_\alpha+j_\alpha+l_\alpha)+\rho_{\alpha-1}=(p-1)+p\rho_\alpha
\quad\left(\begin{gathered}\text{où }\unsure{\alpha-1|\rho_{\alpha-1}|\sigma_\alpha}\\
\text{et }\rho_\alpha\text{ sont}\\0\leq\ \ \leq p-1\end{gathered}\right)
\]

\[
\text{\textcircled{$b'$ bis}}\quad\deleted{\unsure{i_{r+1}}}\quad
\underset{\substack{\shortparallel\\0}}{i_{r+1}}+\underset{\substack{\shortparallel\\0}}{j_{r+1}}+l_{r+1}+\rho_r=(p-1)+p\mathbin{.}0
\]

\deleted{\textcircled{$c'$}} \deleted{\unsure{$i+j$}} Ce qui montre que \deleted{\unsure{lorsque}} l'addition des}

% archive sheet number: 32
